\begin{algorithm}[Factoring a Finite Separable Algebra]\label{alg:factorsep}
Let $A$ be a finite separable algebra over $\F_p$.  This
algorithm either shows that $A$ is a field or finds
a nontrivial idempotent in $A$, i.e., an $\eps\in A$
such that $\eps^2 = \eps$ with $\eps\neq 0$ and $\eps\neq 1$.
\begin{enumerate}
\item The dimension of the kernel $V$ of the map $x\mapsto x^p - x$ is
  equal to $k$.  This is because abstractly we have that $A\ncisom
  A_1\times \cdots \times A_k$, with each $A_i$ a finite field
  extension of $\F_p$.
\item If $k=1$ we are done.  Terminate.
\item Otherwise, choose $\alpha \in V$ with $\alpha \not\in \F_p$.
(Think of $\F_p$ as the diagonal embedding of $\F_p$ in
$A_1\times \cdots \times A_k$).
Compute powers of $\alpha$ and find the minimal polynomial $m(X)$
of $\alpha$.
\item Since $V\ncisom \F_p \times \cdots \times F_p$ ($k$ factors),
the polynomial $m(X)$ is a square-free product of linear factors, that
has degree $>1$ since $\alpha\not\in\F_p$.  Thus we can compute
a splitting $m(X) = m_1(X) \cdot m_2(X)$, where both $m_i(X)$ have
positive degree.
\item Use the Euclidean algorithm in $\F_p[X]$ to find
$U_1(X)$ and $U_2(X)$ such that
$$
  U_1 m_1 + U_2 m_2 = 1.
$$
\item Let $\eps = (U_1 m_1)(\alpha)$.  Then we have
$$
  U_1 m_1 U_1 m_1 + U_2 m_2 U_1 m_1 = U_1 m_1,
$$
so since $(m_1 m_2)(\alpha) = m(\alpha)=01$, we have $\eps^2 = \eps$.
Also, since $\gcd(U_1, m_2) = \gcd(U_2, m_1) = 1$,
we have $\eps\neq 0$ and $\eps \neq 1$.
\end{enumerate}
\end{algorithm}

Given Algorithm~\ref{alg:factorsep}, we compute an idempotent
$\eps \in A$, and observe that
$$
 A \isom \Ker(1 - \eps)  \oplus \Ker(\eps).
$$
Since $(1 - \eps) + \eps = 1$, we see that
$(1 - \eps)v + \eps v = v$, so that the sum
of the two kernels equals $A$.
Also, if $v$ is in the intersection of the two kernels,
then $\eps(v) = 0$ and $(1-\eps)(v) =0$, so
$0 = (1-\eps)(v) = v - \eps(v) = v$, so the sum is direct.

\begin{remark}
  The beginning of \cite[\S6.2.4]{cohen:course_ant} suggests that one
  can just randomly find an $\alpha \in A$ such that $A\isom
  \F_p[x]/(m(x))$ where $m$ is the minimal polynomial of $\alpha$.
  This is usually the case, but is {\em wrong in general}, since there
  need {\em not} be an $\alpha \in A$ such that $A \isom
  \F_p[\alpha]$.  For example, let $p=2$ and $K$ be as in
  Example~\ref{ex:dedekind}.  Then $A\isom \F_2 \times \F_2 \times
  \F_2$, which as a ring is not generated by a single element, since
  there are only 2 distinct linear polynomials over $\F_2[x]$.
\end{remark}

\begin{algorithm}[Factoring a General Prime Ideal]\label{alg:genfac}
Let $K=\Q(a)$ be a number field given
by an algebraic integer~$a$ as a root of its
minimal monic polynomial~$f$ of degree~$n$.
We assume that an order $\O$ has been
given by a basis $w_1,\ldots,w_n$
and that~$\O$ contains $\Z[a]$.
For any prime $p\in\Z$, the following
algorithm computes the set of maximal ideals of~$\O$ that contain~$p$.
\begin{enumerate}
\item{} [Check if easy] If $p\nmid \disc(\Z[a]) / \disc(\O)$ (so
$p\nmid [\O:\Z[a]]$), then using Theorem~\ref{thm:fac1} we
factor~$p\O$.
\item{} [Compute radical]
Let $I$ be the \defn{radical} of $p\O$, which is the ideal of
elements $x\in\O$ such that $x^m\in p\O$
for some positive integer~$m$.  Note that $p\O \subset I$, i.e.,
$I\mid p\O$; also~$I$ is the product
of the primes that divide $p$, without multiplicity.
Using linear algebra over the finite field
$\F_p$, we compute a basis for $I/p\O$ by computing
the abelian subgroup of $\O/p\O$ of all nilpotent
elements.  This computes $I$, since $p\O\subset I$.
\item{} [Compute quotient by radical]
Compute an $\F_p$ basis for
$$
  A = \O/I = (\O/p\O)/(I/p\O).
$$
The second equality comes from the fact that $p\O\subset I$.
Note that $\O/p\O$
is obtained by simply reducing the basis $w_1,\ldots, w_n$ modulo~$p$.
Thus this step entirely involves linear algebra modulo $p$.

\item{} [Decompose quotient] The ring $A$ is isomorphic to
the quotient of $\O$ by a radical ideal,
so it decomposes as a product
$A \isom A_1 \times \cdots \times A_k$ of finite fields.
We find such a decomposition explicitly using Algorithm~\ref{alg:factorsep}.

\item{} [Compute the maximal ideals over $p$] Each maximal ideal
  $\p_i$ lying over~$p$ is the kernel of one of the compositions
$$\O \to A \ncisom A_1 \times \cdots \times A_k \to A_i.$$
\end{enumerate}
\end{algorithm}
Algorithm~\ref{alg:genfac} finds all primes of $\O$ that contain the radical $I$ of
$p\O$.  Every such prime clearly contains $p$, so to see that the
algorithm is correct, we prove that the primes $\p$ of $\O$ that
contain~$p$ also contain~$I$.  If $\p$ is a prime of $\O$ that
contains~$p$, then $p\O \subset \p$.  If $x\in I$ then $x^m\in p\O$
for some $m$, so $x^m\in \p$ which implies that $x\in \p$ by the primality
of $\p$.  Thus $\p$ contains $I$, as required.  Note that we do not find the powers of
primes that divide $p$ in Algorithm~\ref{alg:genfac}; that's left to another
algorithm that we will not discuss in this book.

Algorithm~\ref{alg:genfac} was invented by J.~Buchmann and
H.\thinspace{}W. Lenstra, though their paper seems to have never been
published; however, the algorithm is described in detail in
\cite[\S6.2.5]{cohen:course_ant}.  Incidentally, this chapter is based
on Chapters~4 and~6 of \cite{cohen:course_ant}, which is highly
recommended, and goes into much more detail about these algorithms.


%%[[Add discussion about how to compute a $p$-maximal order here.]]
